\section{Conclusiones Relevantes \label{sec:conclusionesrelevantes}}

Este informe presenta una primera aproximaci\'on al problema de planificaci\'on operacional de la flota de buses el\'ectricos de RED, formulado como un Electric Vehicle Scheduling Problem (E-VSP). Se caracteriz\'o el problema en sus tres decisiones centrales, asignaci\'on de viajes a buses, desplazamientos sin pasajeros y programaci\'on de recargas, y se estableci\'o como objetivo minimizar el costo total de operaci\'on sujeto a la cobertura completa de los viajes y a las restricciones de autonom\'ia, horario y capacidad de carga.

El an\'alisis de datos confirm\'o la magnitud del desaf\'io: un d\'ia laboral concentra 417 rutas y 64.502 expediciones, con dos peaks de demanda (8:00 y 18:00) que superan ampliamente la capacidad de carga simult\'anea de los cinco electroterminales (700 espacios) y que coinciden parcialmente con los bloques tarifarios m\'as caros. Esto valida que el problema exige coordinar conjuntamente las decisiones de transporte y de energ\'ia, y no puede resolverse como una simple asignaci\'on de viajes.

La revisi\'on bibliogr\'afica permiti\'o identificar tres enfoques metodol\'ogicos viables (ALNS, Set Partitioning con Generaci\'on de Columnas y Benders, y descomposici\'on espaciotemporal), cada uno con ventajas y limitaciones distintas en escalabilidad, calidad de soluci\'on y complejidad de implementaci\'on. Se seleccion\'o la descomposici\'on espaciotemporal como metodolog\'ia principal por su capacidad de reducir la escala real del problema a subproblemas manejables (asignaci\'on espacial por electroterminal seguida de ventanas horarias superpuestas), manteniendo la posibilidad de integrarse con la segunda metodolog\'ia en los tramos m\'as complejos.

Esta etapa preliminar permiti\'o sentar las bases conceptuales, de datos y metodol\'ogicas del proyecto. El desaf\'io de las siguientes etapas ser\'a traducir esta propuesta en una implementaci\'on computacional concreta, validarla con los datos reales de RED y evaluar si el enfoque logra un equilibrio adecuado entre tratabilidad computacional y calidad de la soluci\'on obtenida.

\section{Pasos Futuros \label{sec:pasosfuturos}}

Habiendo descrito el problema, analizado los datos disponibles y seleccionado la metodolog\'ia de descomposici\'on espaciotemporal como enfoque principal, el trabajo futuro se estructurar\'a en cinco etapas secuenciales:

\begin{enumerate}
\item \textbf{Formular el modelo por fases.} Se desarrollar\'a la formulaci\'on matem\'atica completa de la metodolog\'ia propuesta, separando expl\'icitamente la Fase 1 (asignaci\'on espacial de viajes a electroterminales) y la Fase 2 (descomposici\'on temporal mediante ventanas superpuestas y horizonte rodante). Esto incluye definir con precisi\'on las variables de decisi\'on, la funci\'on objetivo y el conjunto completo de restricciones (cobertura de viajes, factibilidad horaria, evoluci\'on de bater\'ia, capacidad de carga simult\'anea) para cada subproblema.
\item \textbf{Implementar la modelaci\'on computacional.} Se traducir\'a la formulaci\'on a un modelo ejecutable, utilizando un solver de optimizaci\'on (MILP) para los subproblemas de la Fase 2, junto con las rutinas necesarias para la asignaci\'on espacial de la Fase 1 y la transmisi\'on del estado de cada bus (ubicaci\'on, nivel de bater\'ia) entre ventanas consecutivas.
\item \textbf{Probar con datos de RED.} El modelo se ejecutar\'a sobre el caso base ya construido (d\'ia laboral, 417 rutas, 64.502 expediciones), comenzando con instancias reducidas o electroterminales individuales para verificar el correcto funcionamiento del modelo antes de escalarlo a la red completa.
\item \textbf{Validar y ajustar par\'ametros.} Se evaluar\'a la sensibilidad de la soluci\'on al tama\~no de las ventanas temporales, al criterio de agrupamiento espacial en la Fase 1 y a los costos definidos, ajustando estos par\'ametros en caso de detectar infactibilidades recurrentes o soluciones de baja calidad. En esta etapa tambi\'en se evaluar\'a la alternativa de fusionar los electroterminales Los Espinos y Santa Rosa, y se decidir\'a si es necesario incorporar Set Partitioning con Generaci\'on de Columnas y cortes de Benders en las ventanas horarias de mayor complejidad (horas peak).
\item \textbf{Analizar resultados y presentarlos.} Finalmente, se construir\'an los entregables definidos en la Secci\'on \ref{sec:entregable} (jornadas por bus, itinerario de carga, desglose de costos), se calcular\'an los KPIs definidos en la Secci\'on \ref{sec:kpis} y se compararan los resultados obtenidos contra los del caso base, documentando las conclusiones y limitaciones del enfoque para la presentaci\'on final.
\end{enumerate}

La Tabla \ref{tab:gantt} presenta la Carta Gantt planificada para estas cinco etapas a lo largo de las diez semanas siguientes.

\begin{table}[htb]
\begin{center}
\caption[Carta Gantt de trabajo futuro]{Carta Gantt de trabajo futuro}
\label{tab:gantt}
\resizebox{\textwidth}{!}{%
\begin{tabular}{@{}l *{10}{c}@{}}\toprule
\textbf{Actividad} & \textbf{S1} & \textbf{S2} & \textbf{S3} & \textbf{S4} & \textbf{S5} & \textbf{S6} & \textbf{S7} & \textbf{S8} & \textbf{S9} & \textbf{S10} \\ \midrule
Formular el modelo de optimizaci\'on & \cellcolor{ganttgreen} & \cellcolor{ganttgreen} & & & & & & & & \\
Implementar y validar modelo simplificado & & \cellcolor{ganttgreen} & \cellcolor{ganttgreen} & & & & & & & \\
Incorporar restricciones progresivamente & & & \cellcolor{ganttgreen} & \cellcolor{ganttgreen} & \cellcolor{ganttgreen} & \cellcolor{ganttgreen} & & & & \\
Implementar descomposici\'on temporal & & & & & & \cellcolor{ganttgreen} & \cellcolor{ganttgreen} & & & \\
Integrar metodolog\'ia completa & & & & & & & \cellcolor{ganttgreen} & \cellcolor{ganttgreen} & & \\
Escalar metodolog\'ia para caso RED & & & & & & & & \cellcolor{ganttgreen} & \cellcolor{ganttgreen} & \\
Analizar resultados y ajustar par\'ametros & & & & & & & & & \cellcolor{ganttgreen} & \cellcolor{ganttgreen} \\
\bottomrule
\end{tabular}%
}
\end{center}
\end{table}
